\subsection{常系数 n 阶线性方程}

若方程 (33) 中的系数 \(a_{j}\) 为常数，则相应的矩阵 \(A\) 也是常数；写出 \(e^{rt}\) 是一个解，便得到特征方程，即

\[
r ^ {n} - a _ {1} r ^ {n - 1} - \ldots - a _ {n - 1} r - a _ {n} = 0.\tag{40}
\]

设 \(r_j\) ( \(1 \leqslant j \leqslant q\) ) 是这个方程的互异根，\(n_j\) ( \(1 \leqslant j \leqslant q\) ) 是根 \(r_j \left( \sum_{j=1}^{q} n_j = n \right)\) 的重数。由 IV, p.~188 至 194 的结果，对齐次方程而言，每个根 \(r_j\) 都对应着

\[
\mathrm{D} ^ {n} x - a _ {1} \mathrm{D} ^ {n - 1} x - \dots - a _ {n - 1} \mathrm{D} x - a _ {n} x = 0\tag{41}
\]

一个由 \(n_j\) 个线性无关的积分所成的组

\[
u _ {j k} (t) = e ^ {r _ {j} t} p _ {j k} (t),
\]

其中 \(p_{jk}\) 是次数 \(\leqslant n_{j}-1\) 的（复系数）多项式（ \(1\leqslant k\leqslant n_{j}\) ）；此外，如此得到的 n 个积分 \(u_{jk}\) ( \(1\leqslant j\leqslant q\) , \(1\leqslant k\leqslant n_{j}\) ) 线性无关。由此推出，\(n_{j}\) 个多项式 \(p_{jk}\) ( \(1\leqslant k\leqslant n_{j}\) ) 在 t 的次数 \(\leqslant n_{j}-1\) 的多项式所成的空间中线性无关，从而构成该空间（在 C 上）的一个基，因为后者的维数是 \(n_{j}\) 。换言之：

命题 9. 设 \(r_j\) ( \(1 \leqslant j \leqslant q\) ) 是特征方程 (40) 的互异根，\(n_j\) 是根 \(r_j\) ( \(1 \leqslant j \leqslant q\) ) 的重数。那么 \(n\) 个函数 \(t^k e^{r_j t}\) ( \(1 \leqslant k \leqslant n_j\) , \(1 \leqslant j \leqslant q\) ) 是齐次方程 (41) 的线性无关的积分。

可以直接地证明这个结果如下。由方程 (41) 可知，这个方程的每个积分的 \(n^{th}\) 导数在 R 上可导，由此立即推出（对整数 m \textgreater{} n 用归纳法），(41) 的每个积分都有 m 阶导数，也就是说，在 R 上无穷次可导。设 D 是 R 上无穷次可导的复值函数在 C 上所成的（非拓扑）向量空间；映射 \(x \mapsto Dx\) 是这个空间的一个自同态，而方程 (41) 可以写成

\[
f (\mathrm{D}) x = 0\tag{42}
\]

\[
\text { where } f (\mathrm{D}) = \mathrm{D} ^ {n} - a _ {1} \mathrm{D} ^ {n - 1} - \dots - a _ {n - 1} \mathrm{D} - a _ {n} (\mathrm{Alg}., \mathrm{IV}, \mathrm{p}. 4).
\]

命题 10. 设 \(g\) 与 \(h\) 是两个互素多项式，使得 \(f = gh\) 。(42) 的解所成的子空间是下列两个方程的解子空间的直和：

\[
g (\mathrm{D}) x = 0, \quad h (\mathrm{D}) x = 0.
\]

现在，由贝祖恒等式 (Alg., VII. 2, th. 1)，存在两个多项式 \(p(\mathrm{D})\) 与 \(q(\mathrm{D})\) ，使得 \(p(\mathrm{D})g(\mathrm{D}) + q(\mathrm{D})h(\mathrm{D}) = 1\) 。对 (42) 的每个解 \(x\) ，可以写出 \(x = y + z\) ，其中 \(y = p(\mathrm{D})g(\mathrm{D})x\) 且 \(z = q(\mathrm{D})h(\mathrm{D})x\) ；于是 \(h(\mathrm{D})y = p(\mathrm{D})(f(\mathrm{D})x) = 0\) 且 \(g(\mathrm{D})z = q(\mathrm{D})(f(\mathrm{D})x) = 0\) 。另一方面，若 \(g(\mathrm{D})x = 0\) 与 \(h(\mathrm{D})x = 0\) 同时成立，则可推出

\[
x = p (\mathrm{D}) (g (\mathrm{D}) x) + q (\mathrm{D}) (h (\mathrm{D}) x) = 0,
\]

这就完成了证明。

用前面的记号，可以写出

\[
f (\mathrm{D}) = \prod_ {j = 1} ^ {q} (\mathrm{D} - r _ {j}) ^ {n _ {j}}
\]

而命题 10（对 \(q\) 用归纳法）表明，(42) 的解子空间是下列 \(q\) 个方程

\[
(\mathrm{D} - r _ {j}) ^ {n _ {j}} x = 0 \quad (1 \leqslant j \leqslant q).\tag{43}
\]

现在，对每个复数 r 有

\[
\mathrm{D} (e ^ {r t} x) = e ^ {r t} (\mathrm{D} + r) x\tag{44}
\]

所以 (43) 等价于

\[
\mathrm{D} ^ {n _ {j}} \left(e ^ {- r _ {j} t} x\right) = 0
\]

从而它以函数 \(e^{r_{j}t} p_{j}(t)\) 为解，其中 \(p_{j}\) 遍取次数 \(\leqslant n_{j}-1\) 的多项式所成之集；这样就重新得到 IV, p.~194 的命题 9。

假设齐次方程 (41) 已经解出（也就是说，假设特征根已经求出），我们知道常数变易法可以求出非齐次方程

\[
\mathrm{D} ^ {n} x - a _ {1} \mathrm{D} ^ {n - 1} x - \dots - a _ {n - 1} \mathrm{D} x - a _ {n} x = b (t)\tag{45}
\]

的解，其中 \(b(t)\) 是任意的规则函数 (IV, p.~182)；我们指出，若 \(b(t)\) 在一个区间 \(J\) 上无穷次可导，则 (45) 的所有积分都在 J 上无穷次可导。在特殊情形 \(b(t) = e^{\alpha t} p(t)\) （其中 p 是一个（复系数）多项式，\(\alpha\) 是一个任意复数）下，用下面的方法可以更简单地得到 (45) 的一个积分。令 \(x = e^{\alpha t} y\) ；方程

\[
f (\mathrm{D}) x = e ^ {\alpha t} p (t)
\]

由 (44) 可以写成

\[
f (\alpha + \mathrm{D}) y = p (t)
\]

或者，对多项式 \(f(\mathbf{D})\) 应用泰勒公式，

\[
\frac {f ^ {(n)} (\alpha)}{n !} \mathrm{D} ^ {n} y + \frac {f ^ {(n - 1)} (\alpha)}{(n - 1) !} \mathrm{D} ^ {n - 1} y + \dots + \frac {f ^ {\prime} (\alpha)}{1 !} \mathrm{D} y + f (\alpha) y = p (t).\tag{46}
\]

设 \(m\) 是多项式 \(p(t) = \sum_{k=0}^{m} \lambda_k t^{m-k}\) 的次数；若 \(f(\alpha) \neq 0\) （即 \(\alpha\) 不是特征根），则存在唯一的一个多项式 \(u(t) = \sum_{k=0}^{m} c_k t^{m-k}\) （次数为 \(m\) ）是 (46) 的解，因为系数 \(c_k\) 由线性方程组

\[
\begin{array}{r l} f (\alpha) c _ {k} + \binom {m - k + 1} {1} f ^ {\prime} (\alpha) c _ {k - 1} + \binom {m - k + 2} {2} f ^ {\prime \prime} (\alpha) c _ {k - 2} + \dots \\ & + \binom {m} {k} f ^ {(k)} (\alpha) c _ {0} = \lambda_ {k} \quad (0 \leqslant k \leqslant m) \end{array}
\]

所确定，而这个方程组显然有且仅有一个解。反之，若 \(\alpha\) 是一个特征根，h 是它的重数，则前面的计算表明，存在唯一的一个 m 次多项式，使得 \(D^{h}y = v(t)\) 的每个解都是一个积分；换句话说，(46) 的每个多项式解这时都是 \(m + h\) 次的（“共振”）。

